\documentclass[12pt]{article}
\usepackage{notes-common}
\geometry{margin=0.35in}
\AtBeginDocument{\fontsize{14}{17}\selectfont}

% Fri Oct 30 -- Magnetic force on moving charges (Knight 29.7). Same lesson
% for both tracks. See cyclotron-helix-notes-key.tex for the filled-in
% version -- same structure, wording and diagrams, every blank carries the
% answer. Print layout matches biot-savart-intro-notes.tex: small margins,
% bigger base font.
%
% Sec.2 is Knight Example 29.11 (electron accelerated through 500 V, then
% injected into a uniform field; find the field's strength from the radius
% of the resulting circle). Sec.3 is a new (not-from-the-book) problem:
% split v into components parallel/perpendicular to B to get helical
% motion, then find both the helix radius and the pitch (vertical drift
% per revolution).
\newcommand{\eqline}[2][-0.5mm]{%
  \par\vspace{1mm}\noindent #2\par\vspace{#1}%
  \noindent\textcolor{rulehue}{\rule{\linewidth}{0.5pt}}\par\vspace{4mm}%
}

\begin{document}

\noteshead{Magnetic Force on Moving Charges}{}

\goal{
\begin{itemize}[itemsep=5pt,parsep=2pt,topsep=2pt]
\item See why a charge moving across a uniform $\vec B$ field goes in a
circle, and where the cyclotron radius $r=mv/|q|B$ comes from.
\item Work a full numerical example end to end: accelerate a charge through
a potential difference, then find the field strength from how fast it
circles (Example 29.11).
\item Split $\vec v$ into the piece along $\vec B$ (untouched) and the piece
across it (circles), to see why a charge entering at an angle traces a
\emph{helix}, not a circle.
\item Find both the radius of that helix and its \emph{pitch} --- the
distance the charge drifts along $\vec B$ during one full loop.
\end{itemize}}

%============================================================
\nsec{Why a uniform field makes a charge go in a circle}

A charge $q$ moves with speed $v$ perpendicular to a uniform field $\vec B$.
The magnetic force $\vec F=q\vec v\times\vec B$ is always perpendicular to
$\vec v$ --- it only ever turns the velocity, never speeds it up or slows it
down. A force that's always perpendicular to the motion, with constant
magnitude, is exactly what produces \bl[4.5cm] motion.

Set the magnetic force equal to the centripetal force it's providing:

\[
|q|\,v\,B \ =\ \bl[3.5cm]
\]
\vspace{0.8cm}

Solve for the radius of the circle:

\result[1.6cm]{$r_{\text{cyc}} = $}

\checkpoint{Does a \emph{faster} charge circle in a bigger or smaller
radius? Does a \emph{stronger} field give a bigger or smaller radius? Is
either answer the opposite of what you expected?}

%============================================================
\nsec{Worked example: the radius of cyclotron motion}

\textit{Example 29.11. An electron is accelerated from rest through a
potential difference of $500$ V, then injected into a uniform magnetic
field. Once in the field, it completes half a revolution in $2.0$ ns. What
is the radius of its orbit?}

\begin{center}
\begin{tikzpicture}[x=0.9cm,y=0.9cm]
  \draw[line width=1.6pt] (0,-1.6) -- (0,1.6);
  \draw[line width=1.6pt] (1.0,-1.6) -- (1.0,1.6);
  \node[scale=0.78, anchor=east, xshift=-2pt] at (0,-2.0) {$0$ V};
  \node[scale=0.78, anchor=west, xshift=2pt] at (1.0,-2.0) {$500$ V};
  \foreach \x in {1.8,2.6,3.4}{
    \foreach \y in {-1.2,-0.4,0.4,1.2}{
      \node[scale=0.75, gray] at (\x,\y) {$\times$};}}
  \node[scale=0.78] at (3.95,1.6) {$\vec B$ (into page)};
  \draw[-latex, accent, line width=1pt] (-0.7,0) -- (0.9,0);
  \draw[-latex, accent, line width=1pt] (1.0,0) arc (180:0:0.8);
  \filldraw[accent] (-0.7,0) circle (1.5pt);
  \node[scale=0.75] at (-0.7,0.35) {$e^-$};
\end{tikzpicture}
\end{center}
\figtask{Label the $0$ V and $500$ V plates, mark $\vec B$'s direction, and
sketch the electron's path: straight across the gap, then curving once it
enters the field.}

\textbf{Step 1 --- speed entering the field.} Energy conservation,
$K_f+qV_f=K_i+qV_i$, with $K_i=0$:

\[
\tfrac12 m v_f^2 = e\,V_f \qquad\Longrightarrow\qquad v_f = \
\]
\vspace{1.0cm}

\eqline[0.2cm]{$v_f = $}

\textbf{Step 2 --- turn the half-revolution time into a period and
frequency.} Half a revolution takes $2.0$ ns, so a full revolution (the
period $T$) takes \bl[3cm]. The cyclotron frequency is $f=1/T=$ \bl[3cm].

\textbf{Step 3 --- use the cyclotron frequency to get $B$.} The cyclotron
frequency doesn't depend on $v$ or $r$ at all:

\[
f = \dfrac{|q|B}{2\pi m} \qquad\Longrightarrow\qquad B = \
\]
\vspace{1.0cm}

\eqline[0.2cm]{$B = $}

\textbf{Step 4 --- now find the radius}, using $v_f$ from Step 1 and $B$
from Step 3:

\result[1.6cm]{$r_{\text{cyc}} = \dfrac{mv_f}{eB} = $}

\checkpoint{The review in the book compares this radius to a photo of an
electron beam curving in a magnetic field. Does $8.5$ mm sound like a
reasonable size for an electron's orbit in a lab-scale field, or suspiciously
large/small?}

%============================================================
\nsec{When $\vec v$ isn't perpendicular to $\vec B$: helical motion}

Every example so far has had $\vec v\perp\vec B$ exactly. What if the charge
enters at some other angle $\theta$ to the field?

Split $\vec v$ into two pieces, one along $\vec B$ and one across it:

\[
\vec v = \vec v_\parallel + \vec v_\perp\,, \qquad
v_\parallel = v\cos\theta\,, \qquad v_\perp = v\sin\theta
\]

\begin{center}
\begin{tikzpicture}[x=1cm,y=1cm]
  \draw[-latex, gray, line width=1pt] (-0.3,0) -- (2.4,0) node[right,scale=0.8]{$\vec B$};
  \draw[-latex, accent, line width=1.2pt] (0,0) -- (1.6,1.2) node[above,scale=0.8]{$\vec v$};
  \draw[dashed, gray] (1.6,1.2) -- (1.6,0);
  \draw[-latex, blue!55!black, line width=1pt] (0,0) -- (1.6,0) node[below,scale=0.75]{$v_\parallel$};
  \draw[-latex, red!60!black, line width=1pt] (1.6,0) -- (1.6,1.2) node[right,scale=0.75]{$v_\perp$};
  \draw (0.45,0) arc (0:36.5:0.45);
  \node[scale=0.7] at (0.68,0.16) {$\theta$};
\end{tikzpicture}
\end{center}

\begin{itemize}[itemsep=6pt,parsep=3pt,topsep=3pt]
\item The parallel piece, $\vec v_\parallel$, is \bl[3.5cm] to $\vec B$ ---
so $\vec v_\parallel\times\vec B$ gives \bl[3cm] force at all. Left alone,
it just \bl[5cm].
\item The perpendicular piece, $v_\perp$, is exactly \S1's situation all
over again --- it circles with radius \bl[5cm] (write the \S1 formula, with
$v_\perp$ in place of $v$).
\end{itemize}

Put the two motions together --- a uniform drift along $\vec B$, plus a
circle in the plane across $\vec B$ --- and the path traced out is a
\bl[3.5cm]: a \bl[3.5cm].

\checkpoint{At $\theta=90^\circ$ you're back to \S1's pure circle. What does
the path look like in the other extreme, $\theta=0^\circ$?}

\textit{New problem. A proton ($m=1.67\times10^{-27}$ kg, $q=+e=
1.60\times10^{-19}$ C) enters a region of uniform magnetic field
$B=0.50$ T. Its velocity has magnitude $v=4.0\times10^5$ m/s and points
$30^\circ$ away from $\vec B$. Find (a) the radius of the helix it traces,
and (b) the \emph{pitch} --- the distance it drifts along $\vec B$ during
one full revolution.}

\begin{center}
\begin{tikzpicture}[x=1cm,y=1cm]
  \draw[-latex, gray, dashed, line width=0.9pt] (-0.4,-1.1) -- (6.4,-1.1);
  \draw[-latex, gray, line width=1pt] (-0.4,0) -- (6.4,0) node[right,scale=0.8]{$\vec B$};
  \draw[-latex, gray, dashed, line width=0.9pt] (-0.4,1.1) -- (6.4,1.1);
  \foreach \x in {0,0.9,1.8,2.7,3.6,4.5,5.4}{
    \draw[accent, line width=1.1pt] (\x,0) ellipse (0.42 and 1.05);}
  \draw[<->, gray] (0,-1.5) -- (0.9,-1.5);
  \node[scale=0.75] at (0.45,-1.8) {pitch};
\end{tikzpicture}
\end{center}
\figtask{Mark one full turn of the helix and label the pitch you're solving
for.}

\textbf{Step 1 --- components of $v$ relative to $\vec B$}:

$v_\parallel = v\cos30^\circ = $ \bl[3.5cm] \qquad
$v_\perp = v\sin30^\circ = $ \bl[3.5cm]

\textbf{Step 2 --- radius of the helix} (same formula as \S1, using
$v_\perp$):

\result[1.6cm]{$r_{\text{helix}} = \dfrac{mv_\perp}{qB} = $}

\textbf{Step 3 --- period of one full revolution} (this doesn't depend on
$v_\perp$ at all --- why not? see \S1):

\[
T = \dfrac{2\pi m}{qB} = \
\]
\vspace{1.0cm}

\eqline[0.2cm]{$T = $}

\textbf{Step 4 --- the pitch.} During one period $T$, the parallel motion
just moves at constant velocity $v_\parallel$:

\result[1.6cm]{$\text{pitch} = v_\parallel\,T = $}

\checkpoint{If the proton had instead entered at a \emph{smaller} angle
$\theta$ (closer to parallel with $\vec B$), would the helix get fatter or
skinnier? Would it get more tightly or more loosely wound (bigger or
smaller pitch)?}

\end{document}
